% Generated draft: requires mathematical and source review.
\documentclass[11pt]{article}
\usepackage[margin=25mm]{geometry}
\usepackage{fontspec}
\setmainfont{Latin Modern Roman}
\usepackage{amsmath,amssymb,mathtools,unicode-math}
\setmathfont{Latin Modern Math}
\usepackage{enumitem,needspace,xurl,hyperref}
\setlength{\parindent}{0pt}
\setlength{\parskip}{5pt}
\setlength{\emergencystretch}{3em}
\newcommand{\N}{\mathbb{N}}
\newcommand{\Z}{\mathbb{Z}}
\newcommand{\Q}{\mathbb{Q}}
\newcommand{\R}{\mathbb{R}}
\newcommand{\C}{\mathbb{C}}
\newcommand{\F}{\mathbb{F}}
\newcommand{\A}{\mathbb{A}}
\newcommand{\PP}{\mathbb{P}}
\newcommand{\E}{\mathbb{E}}
\newcommand{\eps}{\varepsilon}
\newcommand{\dd}{\,\mathrm{d}}
\begin{document}
\section*{UnsolvedMath ID: 30002112}
\textbf{Source problem code:} catalogue entry
\subsection{Definitions and mathematical statement}
Let \(m\geq 1\) be an integer and let \(\Omega\subset\R^{m+1}\) be a nonempty smooth bounded domain with boundary \(\Sigma=\partial\Omega\). Equip \(\Sigma\) with the metric induced by the Euclidean metric. Write \(\operatorname{Vol}_m(\Sigma)\) for its hypersurface volume and \(\operatorname{Vol}_{m+1}(\Omega)\) for the volume of the domain.

The supplied record does not define \(I(\Omega)\) or \(\lambda_k(\Sigma)\). The conventional interpretation to be checked against the report is
\[
I(\Omega)=\frac{\operatorname{Vol}_m(\Sigma)}{\operatorname{Vol}_{m+1}(\Omega)^{m/(m+1)}},
\]
with \(\lambda_k(\Sigma)\) an eigenvalue of the nonnegative Laplace--Beltrami operator on the closed hypersurface \(\Sigma\), ordered nondecreasingly and counted with multiplicity. This is an intrinsic spectral problem on \(\Sigma\); it requires no boundary condition on \(\Sigma\). The record does not specify whether the first eigenvalue is indexed by \(0\) or \(1\), whether zero eigenvalues are included, or whether connectedness of \(\Sigma\) is assumed. These conventions remain unresolved.

Subject to confirmation of these definitions and conventions, the proposed assertion is that for every integer \(m\geq1\) there exist finite constants \(A_m,B_m\geq0\), depending only on \(m\), such that for every such domain \(\Omega\) and every positive integer eigenvalue index \(k\),
\[
\lambda_k(\Sigma)\operatorname{Vol}_m(\Sigma)^{2/m}
\leq A_m I(\Omega)^{1+2/m}+B_m k^{2/m}.
\]
The positive-dimensional range and the positive integer range for \(k\) make explicit the intended interpretation of the displayed exponents and spectral index; their precise source conventions require confirmation. The question concerns this additive estimate uniformly over domains and indices.
\subsection{Short English statement}
Does every positive dimension admit constants depending only on that dimension which bound the scale-invariant boundary eigenvalues by an additive combination of an isoperimetric term and \(k^{2/m}\)? The meanings of the isoperimetric ratio, the spectral indexing, and the connectedness assumptions require confirmation from the source. The supplied record labels the problem open; this draft does not independently verify that status.
\subsection{Sources}
\begin{itemize}
\item Local UnsolvedMath registry: supplied complete record for permanent ID \(30002112\), title ``Additive Eigenvalue Bounds for Domain Boundaries,'' collection ``Oberwolfach Reports Open Problems,'' exact source problem code \texttt{catalogue entry}. The record gives proposed year \(2012\) and supplies no named proposer.
\item Supplied report citation: \emph{Geometric Aspects of Spectral Theory (2013)}, Oberwolfach Reports, pp.\ 2013--2076, DOI \texttt{10.4171/owr/2012/33}. \url{https://doi.org/10.4171/owr/2012/33}. The background describes Oberwolfach Reports as documenting mathematical workshops and their open-problem discussions.
\item Supplied literature source, ``OWR report'': \url{https://ems.press/content/serial-article-files/46403}.
\item Supplied literature source, ``Colbois El Soufi Girouard'': \url{https://doi.org/10.1515/crelle-2012-0008}. No article title or fuller bibliographic details are supplied.
\item Literature assessment dated exactly \texttt{2026-08-21}: the supplied curation labels the problem open, states that the report posed this additive estimate as open, describes the cited primary work as proving a different multiplicative isoperimetric bound, and reports that no resolution of the additive form was found. This is supplied dated curation, not an independently verified literature assessment; the background explicitly requests independent verification of time-sensitive status and cited claims.
\item The registry labels the statement \texttt{corrected\_verified} and reports that the official report was checked, that the complete statement was confirmed and preserved, and that previously undelimited mathematical notation was normalized. These are supplied source-verification claims; no source was independently visited for this draft.
\end{itemize}
\end{document}
